\chapter{INTRODUCTION}

\section{Context and Motivation}
In Vietnam, many landlords of boarding houses and small rental buildings still manage their business using simple methods. They often use paper notebooks, \textbf{Excel} files, social chat applications, or check bank transfers manually. Although these methods are easy to use at first, they become very hard to manage when the landlord has more rooms and tenants. Landlords have to spend a lot of time doing things like updating room availability, collecting tenant information, writing contracts, calculating bills every month, tracking payments, answering questions, and fixing broken facilities or handling rules violations.

Using different simple tools at the same time causes many problems. For example, room information can be different between what landlords post online and what they write in their private logs. It is also difficult to find old contracts and payment history. Tenants do not have a place to check their rental information, payment status, and reported issues by themselves. Also, when landlords calculate and write bills by hand, they can easily forget deadlines, make mistakes in calculations, or send payment notifications late.

Therefore, this project aims to build a website called \textbf{Smart Boarding House Management System} to put all these tasks together in one \textbf{SaaS} platform. This system has three main parts: a \textbf{public page} to show available rooms, a \textbf{landlord dashboard}, and a \textbf{tenant portal}. The system also has automatic and smart features like creating bills on schedule, creating \textbf{PDF} invoices with \textbf{VietQR} payment codes, sending notifications, and using \textbf{AI} to chat with users or generate announcements.

\section{Project Overview}
The system has two main user groups: \textbf{landlords} and \textbf{tenants}. The first group is \textbf{landlords}, who can manage rooms, tenants, contracts, invoices, reported incidents, consultation requests, and tenant violations. The second group is \textbf{tenants}, who can view their current contract, track their invoices, report room incidents, receive notifications, and chat with an \textbf{AI} assistant that understands their rental details.

For the \textbf{backend}, \textbf{NestJS} is used to build a modular system. The system is organized into different modules like login, rooms, tenants, contracts, invoices, rent requests, incidents, violations, PDF generation, upload, notification, and AI assistant. The \textbf{frontend} is built using \textbf{Next.js}. It is divided into public pages, login pages, a landlord admin dashboard, and a tenant portal.

\section{Report Structure}
This report is divided into six chapters:

\begin{itemize}
    \item \textbf{Chapter 1 (Introduction):} Introduces the project, its background, reasons for development, and system overview.
    \item \textbf{Chapter 2 (Objectives):} Describes the objectives, functional goals, scope of work, and expected results.
    \item \textbf{Chapter 3 (Requirements Analysis):} Shows the requirements analysis, including user roles, system requirements, use cases, and database details.
    \item \textbf{Chapter 4 (Methodology):} Explains the methodology, such as the chosen technologies, application architecture, database model, and step-by-step implementation.
    \item \textbf{Chapter 5 (Results and Discussion):} Discusses the results, focusing on the implemented features, testing, and limitations.
    \item \textbf{Chapter 6 (Conclusion and Future Work):} Concludes the report and suggests some future improvements.
\end{itemize}

